\chapter{The theorem over number fields}
\label{ch:numberfield}

Theorem~\ref{thm:main-arith} takes as input a complete discrete valuation ring $O$ of
characteristic $0$ with algebraically closed residue field, together with arithmetic
data: a subfield $K$ of $\operatorname{Frac} O$ and a local ring $O_0$ with finite
residue field, mapping into $O$ and into $K$ and covering the elements of $K$ of
valuation $\le 1$. Over $\QQ$ these are $O = W(\Fpbar)$ and $O_0 = \ZZ_{(p)}$; the
theorems over $\QQ$ (Theorem~\ref{thm:main}, Corollary~\ref{cor:finiteness}) are derived
from the number-field statement as the case $A = \ZZ$, $K = \QQ$, $\mathfrak p = (p)$. This chapter supplies
them for an arbitrary number field $K$ and a prime $\mathfrak p$ of its ring of integers; the resulting statement,
Theorem~\ref{thm:main-nf}, is Corollary~6.7 of Stoll, \emph{Independence of rational
points on twists of a given curve}, Compositio Math.\ 142 (2006).

The ring $O$ has to be, up to the choices involved, the ring of integers of the completed
maximal unramified extension of the completion $K_{\mathfrak p}$. Rather than developing
the theory of complete discretely valued fields --- of which Mathlib (as of
version~v4.35.0) has very little: there is no Henselization, no maximal unramified extension, and no structure
theory of complete discrete valuation rings --- we construct $O$ directly, as a direct
factor of the $p$-adically complete ring $W(\Fpbar) \otimes_{\ZZ} O_K$. The merit of this
presentation is that the one genuinely local-arithmetic input --- a uniformizer of
$(O_K)_{\mathfrak p}$ remains a uniformizer of $O$, which is where ``$W(\Fpbar)$ is
unramified over $\Zp$'' enters --- reduces to the separability of the residue extension
$\Fq/\Fp$, in the form of an idempotent splitting of $\Fpbar \otimes_{\Fp} \Fq$.
Everything else is elementary commutative algebra around $p$-adic completeness.

Throughout the chapter, $K$ is a number field of degree $n$ with ring of integers $O_K$;
$p$ is a prime number and $\mathfrak p$ a maximal ideal of $O_K$ containing $p$, with
residue field $\Fq = O_K/\mathfrak p$ and ramification index $e = e(\mathfrak p \mid p)$,
so that $p\,(O_K)_{\mathfrak p} = \mathfrak p^e (O_K)_{\mathfrak p}$; the ideal
$J \le O_K$ denotes the complementary factor, so that $p\,O_K = \mathfrak p^e J$ with
$\mathfrak p + J = O_K$; $\sigma \colon \Fq \to \Fpbar$ is a fixed embedding; and
$W := W(\Fpbar)$ is the ring of Witt vectors, a complete discrete valuation ring with
maximal ideal $pW$ and residue field $\Fpbar$.

\begin{definition}[The base-change ring]
\label{def:witt-tensor}
\lean{ChabautyColeman.NumberFieldBase.rho, ChabautyColeman.NumberFieldBase.extIdeal,
  ChabautyColeman.NumberFieldBase.rho_surjective,
  ChabautyColeman.NumberFieldBase.rho_eq_zero_of_mem_extIdeal}
\leanok
Let $R := W \otimes_{\ZZ} O_K$, a free $W$-module of rank $n$, and let
$\rho \colon R \to \Fpbar$ be the ring homomorphism determined by
$w \otimes a \mapsto \bar w \cdot \sigma(\bar a)$, where $\bar w$ is the residue of $w$
and $\bar a$ the class of $a$ modulo $\mathfrak p$. It is surjective and kills $pR$ and
$\mathfrak p R$. Being module-finite over $W$, the ring $R$ is $p$-adically complete and
Hausdorff, and so is every quotient ring of $R$.
\end{definition}

\begin{lemma}[Separable corners split off]
\label{lem:residue-split}
\lean{exists_isIdempotentElem_ker_eq_span}
\leanok
Let $B$ be a commutative algebra over a field $F$, generated by a single element $t$, and
let $\rho \colon B \to F$ be an $F$-algebra homomorphism. If $t$ satisfies a polynomial
$m \in F[x]$ with $m'(\rho(t)) \neq 0$, then the kernel of $\rho$ is generated by an
idempotent: there is an idempotent $\gamma \in B$ with $\ker \rho = (1 - \gamma)$.

In particular, for a finite separable field extension $k'/k$, a field extension $\Omega$
of $k$ and a $k$-embedding $\sigma \colon k' \to \Omega$, the kernel of the multiplication
map $\mu_\sigma \colon \Omega \otimes_k k' \to \Omega$, $x \otimes a \mapsto x\,\sigma(a)$,
is generated by an idempotent: $\Omega \otimes_k k'$ is generated over $\Omega$ by
$1 \otimes \xi$ for a primitive element $\xi$ of $k'/k$ (for finite $k'$ unconditional:
the unit group is cyclic), which satisfies its separable minimal polynomial, and
separability says exactly that every root is simple.
\end{lemma}
\begin{proof}
\leanok
Write $\alpha := \rho(t)$, so $m(\alpha) = \rho(m(t)) = 0$ and $m = (x - \alpha)\,q$ in
$F[x]$. Since $q(\alpha) = m'(\alpha) \neq 0$, the factors are coprime: there is a
B\'ezout identity $u \cdot (x - \alpha) + v \cdot q = 1$. Set $\gamma := (vq)(t)$. Then
$\gamma^2 - \gamma = (vq \cdot (vq - 1))(t) = -(uv \cdot m)(t) = 0$, so $\gamma$ is an
idempotent, and $1 - \gamma = (u \cdot (x-\alpha))(t)$ lies in the kernel of $\rho$
(evaluate at $\alpha$). Conversely an element of $\ker \rho$ is $g(t)$ with
$g(\alpha) = 0$ (here generation by $t$ enters), so $g = (x - \alpha)h$; and
$(x-\alpha)(t) = \bigl((1 - vq)(x - \alpha)\bigr)(t)$ because the difference is
$v \cdot m$ --- so $g(t) \in (1 - \gamma)$.
\end{proof}

\begin{lemma}[The idempotent]
\label{lem:idempotent}
\lean{ChabautyColeman.NumberFieldBase.exists_isIdempotentElem_ker_rho,
  ChabautyColeman.NumberFieldBase.exists_isIdempotentElem_ker_lift_rho}
\uses{def:witt-tensor, lem:residue-split}
\leanok
There is an idempotent $\varepsilon \in R$ with $\varepsilon \in JR$ whose class modulo
$\mathfrak p R$ is the idempotent $\gamma$ of Lemma~\ref{lem:residue-split}, applied to
the $\Fpbar$-algebra $R/\mathfrak p R$ --- generated by the class $t$ of $1 \otimes \xi$
for a lift $\xi \in O_K$ of a primitive element of $\Fq/\Fp$, which satisfies the
separable minimal polynomial of that element --- and to the homomorphism induced by
$\rho$.
\end{lemma}
\begin{proof}
\leanok
Since $\mathfrak p R + JR = R$, choose $c \in R$ with $c \equiv \gamma \bmod \mathfrak p R$
and $c \in JR$. Its defect $c^2 - c$ lies in
$\mathfrak p R \cap JR = \mathfrak p J R$ (the ideals are coprime). The Newton step
$c \mapsto 3c^2 - 2c^3$ squares the defect ideal --- if $c^2 - c \in I$ then
$(3c^2-2c^3)^2 - (3c^2-2c^3) \in I^2$ --- and moves $c$ only by an element of $I$.
Since $(\mathfrak p J)^{me} \subseteq (\mathfrak p^e J)^m = p^m O_K$ for every $m$, the
iterates form a $p$-adic Cauchy sequence whose defects tend to $0$; by completeness of
$R$ they converge to an idempotent $\varepsilon$ with
$\varepsilon \equiv c \bmod \mathfrak p J R$, which therefore still satisfies
$\varepsilon \equiv \gamma \bmod \mathfrak p R$ and $\varepsilon \in JR$.

\emph{Formalization.} The convergence argument is the existing idempotent-lifting lemma
for adically complete rings (\texttt{IsAdicComplete.exists\_isIdempotentElem\_of\_sq\_sub\_mem}),
applied with the ideal $\mathfrak p J R$, whose adic topology agrees with the $p$-adic
one by the displayed inclusions (the transfer is
\texttt{isAdicComplete\_of\_le\_of\_pow\_le}). The congruence
$\varepsilon \equiv \gamma \bmod \mathfrak p R$ is packaged in the form the next steps
consume: $\ker \rho = \mathfrak p R + (1 - \varepsilon)$, pulled back from
$\ker(\bar\rho) = (1 - \gamma)$ on $R/\mathfrak p R$
(\texttt{exists\_isIdempotentElem\_ker\_lift\_rho}, which applies
Lemma~\ref{lem:residue-split} in an unbundled form to avoid carrying an
$\overline{\mathbb F}_p$-algebra structure on $R/\mathfrak p R$).
\end{proof}

\begin{definition}[The local factor]
\label{def:local-factor}
\lean{ChabautyColeman.NumberFieldBase.Factor, ChabautyColeman.NumberFieldBase.resO,
  ChabautyColeman.NumberFieldBase.rhoO, ChabautyColeman.NumberFieldBase.FactorData,
  ChabautyColeman.NumberFieldBase.map_eq_top_of_mem,
  ChabautyColeman.NumberFieldBase.span_natCast_eq_pow,
  ChabautyColeman.NumberFieldBase.isAdicComplete_map}
\uses{lem:idempotent}
\leanok
Let $O := R/(1 - \varepsilon)$. Since $\rho(\varepsilon) = \mu_\sigma(\gamma) = 1$, the
homomorphism $\rho$ descends to a surjection $\rho \colon O \to \Fpbar$. As a
$W$-module, $O$ is a direct summand of $R$ (the image of $\varepsilon$), hence
torsion-free and module-finite over $W$, and $p$-adically complete and Hausdorff.
Because $\varepsilon \in JR$ maps to $1$, we have $JO = O$, and therefore
\[
  \mathfrak p^e O \;=\; \mathfrak p^e J\, O \;=\; p\,O.
\]
\end{definition}

\begin{lemma}[The factor is a complete discrete valuation ring]
\label{lem:factor-dvr}
\lean{ChabautyColeman.NumberFieldBase.FactorData.isDiscreteValuationRing,
  ChabautyColeman.NumberFieldBase.FactorData.isUniformizer,
  ChabautyColeman.NumberFieldBase.FactorData.exists_unit_mul_pow,
  ChabautyColeman.NumberFieldBase.isLocalRing_factor,
  ChabautyColeman.NumberFieldBase.maximalIdeal_factor,
  ChabautyColeman.NumberFieldBase.map_eq_span,
  ChabautyColeman.NumberFieldBase.hasUniqueDiv_factor,
  ChabautyColeman.NumberFieldBase.charZero_factor,
  ChabautyColeman.NumberFieldBase.isAlgClosed_residueField,
  ChabautyColeman.NumberFieldBase.isAdic_map,
  ChabautyColeman.NumberFieldBase.t2Space_factor,
  ChabautyColeman.NumberFieldBase.completeSpace_factor}
\uses{def:local-factor, lem:residue-split}
\leanok
Choose $\varpi \in \mathfrak p \setminus \mathfrak p^2$. The ring $O$ is a complete
discrete valuation ring of characteristic $0$: it is local with maximal ideal
\[
  \mathfrak m \;=\; \ker \rho \;=\; \mathfrak p\, O \;=\; \varpi\, O,
\]
every element of $O_K \setminus \mathfrak p$ is a unit of $O$, every nonzero $x \in O$ is
$\varpi^a u$ with $u \in O^\times$ and $a = v(x)$ for the normalized valuation $v$, the
residue field is $\Fpbar$ (via $\rho$), and $v(p) = e$; in particular $p \mid \varpi^e$,
and the $\mathfrak m$-adic topology is the $p$-adic one, for which $O$ is complete and
Hausdorff.
\end{lemma}
\begin{proof}
\leanok
Modulo $\mathfrak p$, the ring $O/\mathfrak p O \cong
(\Fpbar \otimes_{\Fp} \Fq)/(1-\gamma)$ is the corner cut out by $\gamma$, which
Lemma~\ref{lem:residue-split} identifies with $\Fpbar$ via $\mu_\sigma$; so
$\mathfrak p O = \ker \rho$ is a maximal ideal. For $x \notin \mathfrak p O$ pick $y$
with $xy \equiv 1 \bmod \mathfrak p O$; then $xy = 1 - t$ with $t \in \mathfrak p O$, and
$\sum_j t^j$ converges $p$-adically because $t^{me} \in \mathfrak p^{me} O = p^m O$; so
$x$ is a unit, and $O$ is local with $\mathfrak m = \mathfrak p O$ and residue field
$\Fpbar$. An element $s \in O_K \setminus \mathfrak p$ has $\rho(s) = \sigma(\bar s)
\neq 0$, so $s \notin \mathfrak m$ is a unit; hence $O$ receives $(O_K)_{\mathfrak p}$
and $\mathfrak p O = \mathfrak p (O_K)_{\mathfrak p} O = \varpi O$.

The two filtrations $(\mathfrak m^a)_a$ and $(p^m O)_m$ are cofinal by
$\mathfrak m^e = pO$, so $O$ is $\mathfrak m$-adically complete and Hausdorff, i.e.\
$\bigcap_a \varpi^a O = 0$. Since $O$ is a nonzero direct summand of a free $W$-module,
$p$ is regular on $O$ and no $p^m$ vanishes; from $pO = \varpi^e O$ we get
$p = \varpi^e w$ and $\varpi^e = p w'$ with $w w' = 1$ (cancel the regular $p$), so
$\varpi^{em} = p^m \cdot \mathrm{unit} \neq 0$ and no power of $\varpi$ vanishes. Now
every nonzero $x \in O$ lies in some largest $\varpi^a O$ (Hausdorffness), so
$x = \varpi^a u$ with $u \notin \varpi O = \mathfrak m$ a unit. Products of nonzero
elements are then of the form $\varpi^{a+b}\,\mathrm{unit} \neq 0$, so $O$ is a domain;
a nonzero ideal is generated by $\varpi^a$ for the minimal occurring $a$, so $O$ is a
principal ideal domain, and a discrete valuation ring with $v(x) = a$, $v(p) = e$ and
$p \mid \varpi^e$. Finally $O$ is torsion-free over $\ZZ$, so of characteristic $0$.

\emph{Formalization.} The factor is a definition (\texttt{Factor}) rather than the raw
quotient, so that it carries one canonical ring structure --- the raw quotient of the
tensor ring has several definitionally equal instance paths, which block rewriting. The
uniformizer is certified by an ideal identity $(\varpi) = \mathfrak p\,C$ with
$C \not\subseteq \mathfrak p$ (equivalent to $\varpi \in \mathfrak p \setminus
\mathfrak p^2$, but directly usable: for $x \in \mathfrak p$ and $s \in C \setminus
\mathfrak p$, the product $xs$ lies in $(\varpi)$ and $s$ is a unit of $O$); all the
certificates are bundled as \texttt{FactorData}. Torsion-freeness is proved on
coordinates for a base-changed basis of $R$, and descends to $O$ by multiplying with
$\varepsilon$; the discrete-valuation structure enters Mathlib through
\texttt{IsDiscreteValuationRing.ofHasUnitMulPowIrreducibleFactorization}.
\end{proof}

\begin{lemma}[The arithmetic data]
\label{lem:factor-arithmetic}
\uses{lem:factor-dvr}
\lean{ChabautyColeman.NumberFieldBase.locFactor,
  ChabautyColeman.NumberFieldBase.isLocalHom_locFactor,
  ChabautyColeman.NumberFieldBase.comp_algebraMap_eq_comp_locFactor,
  ChabautyColeman.NumberFieldBase.resO_injective,
  ChabautyColeman.NumberFieldBase.exists_algebraMap_eq,
  ChabautyColeman.NumberFieldBase.residueFieldEquiv,
  ChabautyColeman.NumberFieldBase.finite_residueField,
  ChabautyColeman.NumberFieldBase.finite_quotient}
\leanok
Set $O_0 := (O_K)_{\mathfrak p}$, a discrete valuation ring with finite residue field
$\Fq$.
\begin{enumerate}
\item The map $O_K \to O$ extends to a local homomorphism $\iota \colon O_0 \to O$.
\item $O_K \to O$ is injective, so $K = \operatorname{Frac} O_K$ embeds into
  $\operatorname{Frac} O$, compatibly with $\iota$ and with the inclusion
  $\iota_0 \colon O_0 \to K$.
\item The valuation of $\operatorname{Frac} O$ restricts on $K$ to the $\mathfrak p$-adic
  valuation: $v(x) = v_{\mathfrak p}(x)$ for $x \in K^\times$. Consequently every element
  of $K$ that is integral for the valuation of $O$ --- the covering condition ``of
  valuation $\le 1$'' of Theorem~\ref{thm:main-arith} --- lies in $O_0$.
\end{enumerate}
\end{lemma}
\begin{proof}
\leanok
(1) is Lemma~\ref{lem:factor-dvr}: the elements of $O_K \setminus \mathfrak p$ become
units of $O$, and $\iota$ is local because $\mathfrak p O_0$ maps into
$\mathfrak p O = \mathfrak m$. For (2), the kernel of $O_K \to O$ is a prime ideal; it is
not maximal, since a maximal kernel $\mathfrak q$ would put the positive characteristic
$\ell$ of $O_K/\mathfrak q \hookrightarrow O$ into the kernel of $\ZZ \to O$,
contradicting characteristic $0$; and $O_K$ has Krull dimension $1$, so the kernel is
$0$. For (3), write $x \in K^\times$ as $\varpi^m u$ with $m = v_{\mathfrak p}(x)$ and
$u \in O_0^\times$; then $\iota(u)$ is a unit of $O$ and $v(x) = m$. The consequence
also follows directly from $O_0$ being a valuation ring of $K$: if $x \notin O_0$, then
$x^{-1}$ is a non-unit of $O_0$, so $\iota(x^{-1})$ is a non-unit of $O$; but if $x$ is
integral for the valuation of $O$, then $x = z$ for some $z \in O$, and
$z\,\iota(x^{-1}) = 1$ exhibits $\iota(x^{-1})$ as a unit.

\emph{Formalization.} $O_0$ is Mathlib's valuation subring of $K$ at $\mathfrak p$
(\texttt{valuationSubringAtPrime}), so $\iota_0$ is the inclusion, and the second
argument for (3) is the one formalized; the residue field of $O_0$ is identified with
$O_K/\mathfrak p$ by \texttt{residueFieldEquiv}.
\end{proof}

\begin{theorem}[Chabauty--Coleman over number fields]
\label{thm:main-nf}
\uses{thm:main-arith, lem:factor-dvr, lem:factor-arithmetic}
\lean{ChabautyColeman.NumberFieldBase.main_affine_numberField,
  ChabautyColeman.NumberFieldBase.main_affine_numberField_int,
  ChabautyColeman.NumberFieldBase.separable_of_squarefree_map_residue,
  ChabautyColeman.NumberFieldBase.ringChar_quotient_prime,
  ChabautyColeman.NumberFieldBase.under_eq_span_ringChar,
  ChabautyColeman.NumberFieldBase.ringChar_residueField,
  ChabautyColeman.NumberFieldBase.nonempty_ringHom_algebraicClosure,
  ChabautyColeman.NumberFieldBase.exists_pow_mul,
  Ideal.exists_span_singleton_eq_mul_of_ne_bot,
  ChabautyColeman.NumberFieldBase.squarefree_map_locFactor,
  ChabautyColeman.NumberFieldBase.encard_residuePairs_nf}
\leanok
Let $K$ be a number field and $C \colon y^2 = f(x)$ a hyperelliptic curve of genus $g$ over
$K$ with $f$ monic of odd degree $2g+1$, and let $\mathfrak p$ be a prime of $O_K$ with residue field $\Fq$, residue characteristic $p$
and ramification index $e$, such that $f$ has $\mathfrak p$-integral coefficients and the
reduction $\bar f \in \Fq[x]$ is squarefree.
If $\operatorname{rk} J(K) \le r$ with $r < g$ and
\[
  p \;>\; 2r + e + 1,
\]
then
\[
  \#C(K) \;\le\; \#\bar C(\Fq) + 2r.
\]
This is Corollary~6.7 of Stoll, \emph{loc.\ cit.}
\end{theorem}
\begin{proof}
\leanok
Instantiate Theorem~\ref{thm:main-arith} at the ring $O$ of
Definition~\ref{def:local-factor}, a complete discrete valuation ring of characteristic
$0$ with algebraically closed residue field $\Fpbar$ (Lemma~\ref{lem:factor-dvr}), with
uniformizer $\varpi$ and $v(p) = e$, and at the arithmetic data $K$,
$O_0 = (O_K)_{\mathfrak p}$ of Lemma~\ref{lem:factor-arithmetic}, whose part (3) is the
required covering of $\{x \in K : v(x) \le 1\}$ by $O_0$. The polynomial $f$ has
coefficients in $O_0$, so it is the image of a monic polynomial over $O_0$, whose reduction
modulo the maximal ideal is $\bar f$; the theorem is applied to that polynomial, and the
residue field of $O_0$ is $\Fq$, so the point count on the right-hand side is
$\#\bar C(\Fq)$. The residue characteristic $p$ is odd, indeed $p > 2r + e + 1 \ge 2$
because $e \ge 1$. The reduction of $f$ over the residue field of $O$ is the image of
$\bar f$ under $\sigma \colon \Fq \to \Fpbar$, and squarefreeness persists: over the
perfect field $\Fq$ a squarefree polynomial is separable, and separability is preserved
by any field extension.

\emph{Formalization.} The Lean statements of this theorem and of
Corollary~\ref{cor:finiteness-nf} are phrased over an arbitrary Dedekind domain $A$,
finite free over $\ZZ$, with fraction field $K$ --- so that $A = \ZZ$ itself qualifies at
$K = \QQ$, and the statements over $\QQ$ become direct instantiations. The prime is a
point $v$ of the height-one spectrum of $A$ (a finite place of $K$), $O_0$ is Mathlib's
ring of $v$-integers of $K$ (\texttt{valuationSubringAtPrime}), and the
$\mathfrak p$-integrality of $f$ is certified by a polynomial over $O_0$ mapping to $f$,
whose reduction modulo the maximal ideal is $\bar f$ and whose points over the residue
field are counted. No auxiliary data appears: the residue characteristic is the
characteristic of that residue field (\texttt{ringChar}), and the ramification index is
taken at the ideal of $\ZZ$ below $\mathfrak p$ (\texttt{Ideal.ramificationIdx'
($\mathfrak p$.under $\ZZ$) $\mathfrak p$}). The version for $f \in A[x]$, with the
reduction taken in $A/\mathfrak p$ (\texttt{main\_affine\_numberField\_int}), is a
corollary.
\end{proof}

\begin{corollary}[Finiteness over number fields]
\label{cor:finiteness-nf}
\uses{thm:main-nf}
\lean{ChabautyColeman.NumberFieldBase.ratPoints_finite_numberField,
  ChabautyColeman.NumberFieldBase.ratPoints_finite_numberField_int,
  ChabautyColeman.NumberFieldBase.exists_prime_squarefree_reduction,
  ChabautyColeman.NumberFieldBase.integralModel,
  ChabautyColeman.NumberFieldBase.dvd_absNorm,
  ChabautyColeman.NumberFieldBase.ramificationIdx'_le_finrank,
  Polynomial.separable_of_isUnit_resultant}
\leanok
Let $K$ be a number field and $C \colon y^2 = f(x)$ a hyperelliptic curve of genus $g$ over
$K$ with $f$ monic of odd degree $2g+1$ (squarefreeness of $f$ is part of being a curve). If
$\operatorname{rk} J(K) < g$, then $C(K)$ is finite. No prime appears in the
statement: a suitable prime exists for every such $f$.
\end{corollary}
\begin{proof}
\leanok
Choose $b \in O_K \setminus \{0\}$ with $g := b f \in O_K[x]$, and let $R \in K$ be the
resultant of $f$ and $f'$; it is nonzero, since $f$ is squarefree, hence separable, over
the perfect field $K$. Write $R = a/s$ with $a, s \in O_K$ nonzero. Let $r < g$ bound the
rank and choose a rational prime $p$ beyond both $2r + [K:\QQ] + 1$ and the absolute norm
$N\bigl((asb)\bigr)$, and a prime $\mathfrak p$ of $O_K$ above $p$ (one exists: $p$
divides $N\bigl((p)\bigr) = p^{[K:\QQ]}$). Then $asb \notin \mathfrak p$ --- otherwise
$p \mid N(\mathfrak p) \mid N\bigl((asb)\bigr)$, where $p \mid N(\mathfrak p)$ because
$O_K/\mathfrak p$ is a finite field of characteristic $p$ --- so none of $a$, $s$, $b$
lies in $\mathfrak p$. Hence $f = b^{-1} g$ has $\mathfrak p$-integral coefficients, and
the resultant of this integral model with its derivative is $a/s$, a unit of
$(O_K)_{\mathfrak p}$; its residue is nonzero, so the resultant of $\bar f$ and
$\bar f'$ over $\Fq$ is a unit and $\bar f$ is squarefree. Moreover the ramification
index satisfies $e \le [K:\QQ]$: taking norms in $p\,O_K = \mathfrak p^e J$ gives
$p^{[K:\QQ]} = N(\mathfrak p)^e\,N(J)$ with $p \mid N(\mathfrak p)$. Hence
$p > 2r + e + 1$, and Theorem~\ref{thm:main-nf} bounds $\#C(K)$ by the finite quantity
$\#\bar C(\Fq) + 2r$.

\emph{Formalization.} The version for $f \in O_K[x]$
(\texttt{ratPoints\_finite\_numberField\_int}) is a corollary: $f$ stays squarefree over
$K$ by Gauss's lemma.
\end{proof}
